\FloatBarrier
\section{Recommendations}\label{sec:conclusion}

For a test of the mean under long memory we recommend the following.
\begin{enumerate}
\item Use critical values for the value of $H$ at hand. Short-memory values are invalid under long memory: at $H=0.9$
the nominal $5\,\%$ tests reject a true null between \pthree{naive-min-0.90} and \pthree{naive-max-0.90}\,\% of the time
(\cref{fig:headline}). In the three applications, changing the cutoff changes some decisions and widens nominal
confidence intervals by up to \pthree{emp-nile-HI16-ratio} times (\cref{sec:empirical}).
\item Take critical values from Supplementary Material Tables~S.1--S.2 or the accompanying archive, or, for
fixed-smoothing tests, from the closed-form rule \eqref{eq:rescaled}, whose size error is at most
\pthree{rs-maxerr-090-0.05}\,pp at $5\,\%$ for the tabulated $H\le0.9$ and every tabulated $m\le16$ and $B\le32$. For the
Bartlett/self-normalised test use \cref{tab:sncv} or the response surface checked at held-out grid points of \cref{tab:surface}; the rescaling
rule does not apply to it.
\item Prefer a fixed-smoothing test with few components, $m\le4$ frequencies or $B\in\{4,8\}$ cosines: with the
correct critical value it is within \pthree{fn-maxdev-small-100}\,pp of $5\,\%$ from $n=100$ in every design studied. With
$m=4$ ($B=8$) its limit power exceeds that of the Bartlett/self-normalised test at every $H$ and local alternative of
our grid (\cref{sec:power}); use more components only when $n$ is large or short-run dynamics are known to be weak.
\item Account explicitly for uncertainty in $\hat H$. With a raw local Whittle estimate, nominal $5\,\%$ tests reject
\pthree{pl-pure-min-100}--\pthree{pl-pure-max-100}\,\% of the time at $n=100$ even without short-run dynamics
(\cref{tab:plugin}), while the critical-value error itself is negligible (below \pthree{sn-print-maxerr}\,pp for the printed Bartlett values). For HI with $m=4$, ELW brought size closer to nominal in \pthree{comp-elw-closer}
of the 24 designs. The upper-confidence rule reduced both size and power. Even with that rule, the simulated
rejection rate reached \pthree{comp-up-worst}\,\% at $n=100$ (\cref{sec:competitors}).
\item When comparing with earlier tables, match the normalisation and tail of the statistic (\cref{tab:norm}) first.
Published simulation tables deliver Monte Carlo precision of a few tenths of a percentage point in size; for strongly
antipersistent series ($\delta\le-0.4$) the fixed-$m$ values of \citet{HI17} reflect an inferred finite design rather
than the limit, and practitioners should compute exact finite-$n$ laws directly via \cref{prop:exactn}.
\end{enumerate}
\CosCert{} extends to other functionals of the same cosine coordinates, such as the stationarity statistics of
\citet{KPSS92} under long memory; at $H=\frac12$ it already reproduces the Cram\'er--von Mises law of $\eta_\mu$
(Supplementary Material Table~S.9). We leave these extensions to future work.
